\documentclass[11pt]{article}
\usepackage[margin=1in]{geometry}
\usepackage{amsmath,amssymb,amsthm}
\usepackage{booktabs}
\usepackage[hyphens]{url}
\usepackage{hyperref}

\newtheorem{theorem}{Theorem}[section]
\newtheorem{metatheorem}[theorem]{Meta-Theorem}
\newtheorem{corollary}[theorem]{Corollary}

\title{\textbf{Breaking Out of the Adelic Semigroup in Mahler's $3/2$ Problem:\\
Post-Bit-Length Obstructions, Ceiling-Dual Gauge Renormalization, Parry Kneading Dynamics, and a 154-Theorem Lean~4 Library}}
\author{Joshua Coon\\
\small Cordoned-Off Artifact: \nolinkurl{semigroup_breakout/SemigroupBreakoutTheorems.lean} (154 Theorems \texttt{SB.1}--\texttt{SB.154}, Lean 4.19.0, 0 \texttt{sorry}s)}
\date{November 2026}

\begin{document}
\maketitle

\begin{abstract}
Let $Z := \{\alpha \in \mathbb{R}_{>0} : \forall n \in \mathbb{N},\, \{\alpha(3/2)^n\} \in [0, 1/2)\}$ denote Mahler's set of $Z$-numbers, conjectured since Mahler (1968)~\cite{Mahler1968} to be empty ($Z = \emptyset$). Building on our 226-file bottom-up Lean~4 formalization campaign, we examine whether formal-kernel-guided proof search merely interpolates inside the \emph{convex hull} of known theorems or explores the \emph{semigroup} generated by elementary algebraic, cocycle, and modular operations. We formulate an \textbf{Adelic Finite-Horizon Decoupling Principle} (Meta-Theorem~\ref{thm:meta-decoupling} and \texttt{SB.109}) explaining why purely local congruence and cocycle chains hit a structural wall: at any finite horizon $K$, the integer state equations $2^{T_r} P_r = 3^{T_r} P_0 + A_{0,r}(s)$ and real companion equations $2^{T_r} h_r = 3^{T_r} h_0 - A_{0,r}(s)$ decouple $P_0$ from $h_0$ and admit simultaneous solutions modulo $2^{T_K+1}$ and modulo every odd modulus $2m+1$, so every Parry-admissible $K$-block word is realized by a positive odd integer $b_0 \in (0, 2^{T_K+1}) \subset \mathbb{N}$.

Guided by the \textbf{Post-Bit-Length Exhaustion Principle} ($u_K \ge 1$ and $P_0 < 2^{T_K}$), we formalize \textbf{154 \texttt{sorry}-free, warning-free theorems (\texttt{SB.1}--\texttt{SB.154})} in \nolinkurl{SemigroupBreakoutTheorems.lean} (\texttt{SB.1}--\texttt{SB.152} over $\mathbb{N}/\mathbb{Z}$ and \texttt{SB.153}--\texttt{SB.154} over $\mathbb{R}$; complex character sums $\mathcal{S}(h;N) \in \mathbb{C}$ and external transcendence bounds are treated analytically). Principal contributions include: (1) formalization of Flatto's (1992)~\cite{Flatto1992} and Flatto--Lagarias--Pollington's (1995)~\cite{FLP1995} effective bound $\#(Z \cap [0, X)) \le |\mathcal{C}_W| \le 3(3/2)^W - 2$ ($X \le 2^W$, $O(X^{\log_2(3/2)}) \approx O(X^{0.58496})$) directly over $\mathbb{R}$ (\texttt{SB.112}, \texttt{SB.153}), the exact Parry kneading spine renewal identity on $|\mathcal{C}_N|$ (\texttt{SB.141}, supporting $O(N^2)$ dynamic programming out to $N = 300$), and exact $(N+1)$ scaled-width cardinality (\texttt{SB.149}); (2) the Ceiling-Dual Gauge $W_r = (5 P_r + 4)/36$ (\texttt{SB.81}), Two-Superblock Catalan--Pillai reduction (\texttt{SB.115}), and unconditional real-companion $159/160$ ($99.375\%$) mod-$320$ even-gauge pinch (\texttt{SB.138}--\texttt{SB.139}, \texttt{SB.154}); (3) Kempner Threshold Inversion proving $3/2$ is uniquely over-determined among odd fractions $q/2$ (\texttt{SB.105}--\texttt{SB.108}); and (4) an analysis of seven structural barriers guarding $Z = \emptyset$.
\end{abstract}

\section{Convex Hull vs.\ Semigroup, Scope of Formalization, and Finite-Horizon Decoupling}

All 154 theorems in \nolinkurl{SemigroupBreakoutTheorems.lean} compile under Lean~4.19.0 (\texttt{-DwarningAsError=true}) using only Lean's foundational axioms (\texttt{propext}, \texttt{Classical.choice}, \texttt{Quot.sound}), with zero \texttt{sorry}, \texttt{admit}, or \texttt{native\_decide}. Statements involving $\mathbb{C}$-valued Fourier transforms, Plancherel unitary steps, operator-rank interpretations, and external Diophantine theorems~\cite{Ridout1957,Polya1921,Schlickewei1990,Matveev2000} are presented analytically in the text.

\begin{metatheorem}[Finite-Horizon Decoupling Principle]\label{thm:meta-decoupling}
At any fixed horizon $K \in \mathbb{N}$ and for any block sequence $s = (s_0, \dots, s_{K-1})$ with $s_r \ge 1$:
\begin{enumerate}
  \item Every integer state $P_r = 3 b_r + 1$ ($1 \le r \le K$) satisfies $2^{T_r} P_r = 3^{T_r} P_0 + A_{0,r}(s)$.
  \item Every real tail $h_r = 3 \alpha (3/2)^{T_r} - P_r$ ($0 \le r \le K$) satisfies $2^{T_r} h_r = 3^{T_r} h_0 - A_{0,r}(s)$.
  \item The variables $P_0 \in \mathbb{N}$ and $h_0 \in \mathbb{R}$ decouple at horizon $K$: $0 < h_r < 1/2$ ($0 \le r \le K$) constrains only $h_0 \in (A_{0,K}/3^{T_K}, (A_{0,K}+2^{T_K-1})/3^{T_K})$, while the integer recurrence constrains only $P_0 \bmod 2^{T_K+1}$.
\end{enumerate}
Hence every Parry-admissible $K$-block word $s$ is realized by a positive odd integer $b_0 \in (0, 2^{T_K+1}) \subset \mathbb{N}$, and no auxiliary polynomial $Z = F(P_0, \dots, P_K)$ can violate $|Z|_\infty \ge 2^a 3^b$ without invoking $P_0 < 2^{T_K}$.
\end{metatheorem}

Empirically (\nolinkurl{loop1_product_formula_and_pade.py}), across horizons $N = 1, \dots, 25$ (up to $124$-bit seeds), the minimal engineered seed $b_0^{(N)} \in (0, 2^{T_N+1})$ survives an average of $2.48$ extra binary steps past $T_N$, dying on the first un-engineered block in $23/25$ horizons. Across $[1, 20{,}000]$ (\nolinkurl{loop3_archimedean_ceiling_and_small_u.py}), all $20{,}000$ seeds escape by step $39$.

\section{Formalized $\mathbb{Z}_2$-Symmetry-Breaking Theorems}

\begin{theorem}[Global Archimedean Rung Ceiling (SB.1)]
\emph{Lean identifier:} \nolinkurl{archimedean_rung_ceiling_and_stepSum_bound}.
For any positive seed $b_0 \in \mathbb{N}$ realizing \texttt{ValidBlockChain b\_0 s u (K + 1)} with $0 < \Sigma_K(s)$:
\[
  1 \le u_K,\qquad 2^{s_K} \le P_K,\qquad 2 \cdot 2^{T_{K+1}} \le 2 \cdot 2^{T_{K+1}} u_K < 3^{T_K}(2 P_0 + 1),\qquad 2^{s_0} \le 3 b_0 + 1.
\]
\end{theorem}

\begin{theorem}[Pure Power-of-3 Collapse \& Mod-64 Trap (SB.2)]
\emph{Lean identifier:} \nolinkurl{unit_quotient_even_rung_instant_death_and_odd_rung_mod64_identity}.
When $u_K = 1$ ($b_{K+1} = 3^{s_K - 1}$):
\begin{enumerate}
  \item If $s_K = 2j$ ($j \ge 1$) is even, $P_{K+1} = 3^{2j} + 1 \equiv 2 \pmod 8$, forcing $s_{K+1} = 1$ (instant death).
  \item If $s_K = 2j + 1$ ($j \ge 1$) is odd, $P_{K+1} = 3^{2j+1} + 1 \equiv 4 \pmod 8$ ($s_{K+1} = 2$), and the next state $b_{K+2} = 3(P_{K+1}/4)$ satisfies
  \[
    4 P_{K+2} = 3^{2j+3} + 13 = 243(3^{2j-2} - 1) + 256,
  \]
  so $64 \mid P_{K+2} \iff 256 \mid (3^{2j-2} - 1)$ ($s_K \equiv 3 \pmod{64}$). At $j = 1$ ($s_K = 3, b_{K+1} = 9$), $b_{K+2} = 21$ has $P_{K+2} = 64 = 2^6 \cdot 1$ ($s_{K+2} = 6$ even, $u_{K+2} = 1$), dying at $b_{K+3} = 243$ ($P_{K+3} = 730 \equiv 2 \pmod 8$).
\end{enumerate}
\end{theorem}

\begin{theorem}[Cross-Window Hankel Annihilation (SB.3)]
\emph{Lean identifier:} \nolinkurl{equal_step_window_hankel_determinant_identity_and_bound}.
For any two block windows with $2^d P_2 = 3^d P_1 + A_{12}$ and $2^d P_4 = 3^d P_3 + A_{34}$:
\[
  2^d(P_1 P_4 - P_2 P_3) = P_1 A_{34} - P_3 A_{12},
\]
and when $2 A_{12} < 3^d$ and $2 A_{34} < 3^d$, $-3^d P_3 < 2^{d+1}(P_1 P_4 - P_2 P_3) < 3^d P_1$.
\end{theorem}

\begin{theorem}[Two-Seed First-Divergence Valuation Law (SB.4)]
\emph{Lean identifier:} \nolinkurl{two_seed_first_divergence_valuation_and_bound}.
For any $b_1 \le b_2$ whose binary parities agree for $k < n_0$ and diverge at step $n_0$:
\[
  2^{n_0}\big(\operatorname{orbit}(b_2, n_0) - \operatorname{orbit}(b_1, n_0)\big) = 3^{n_0}(b_2 - b_1),\qquad \big(\operatorname{orbit}(b_2, n_0) - \operatorname{orbit}(b_1, n_0)\big) \bmod 2 = 1,
\]
so $2^{n_0} \mid (b_2 - b_1)$ and $2^{n_0+1} \nmid (b_2 - b_1)$ ($n_0 = v_2(b_2 - b_1)$). Thus $b_2 - b_1 < 2^N$ forces divergence before step $N$.
\end{theorem}

\section{The Three-Wildcard Expedition: Density-Zero Extinction, Base-3 Frozen Prefixes, and Mahler--Antihydra $BB(6)$ Duality}

\begin{theorem}[4-Lift Double-Odd Exclusion (SB.6)]
\emph{Lean identifier:} \nolinkurl{dyadic_four_lift_double_odd_bijection_and_three_of_four_survival}.
For any prefix seed $r$, step $k$, and 2-bit lift $t$, if parities match on $[0, k)$, then $\operatorname{orbit}(r + 2^k t, k) = \operatorname{orbit}(r, k) + 3^k t$, and $(\varepsilon_k = 1 \wedge \varepsilon_{k+1} = 1) \iff \operatorname{orbit}(r, k) \bmod 4 = 3$. Since $3^k \bmod 4 \in \{1, 3\}$, exactly $1$ of the $4$ lifts $t \in \{0, 1, 2, 3\}$ hits the forbidden double-odd residue $3 \pmod 4$, leaving at most $3^m$ seeds in $[0, 4^m)$ avoiding even-step double-odd transitions up to step $2m$.
\end{theorem}

\begin{theorem}[Base-3 Frozen Top Digits \& Forbidden Trit 2 (SB.7)]
\emph{Lean identifier:} \nolinkurl{raw_orbit_base3_frozen_prefix_and_leading_trit_le_one}.
Whenever $1 \le N$, $2^N g_N = 3^N g_0 + K_N$, and $2 K_N < 3^N$:
\[
  \lfloor 2^N g_N / 3^N \rfloor = g_0,\qquad (2^N g_N) \bmod 3^N = K_N,\qquad \lfloor ((2^N g_N) \bmod 3^N) / 3^{N-1} \rfloor \le 1,
\]
and if $g_0 < 3^L$, then $2^N g_N < 3^{N + L}$.
\end{theorem}

\begin{theorem}[Mahler--Antihydra $BB(6)$ Complement Duality (SB.8)]
\emph{Lean identifier:} \nolinkurl{mahler_antihydra_bb6_complement_duality}.
Let $\operatorname{antihydraOrbit}(a_0, k)$ be the $k$-th iterate of Antihydra's floor map $a \mapsto \lfloor 3a/2 \rfloor$. For every $N \in \mathbb{N}$ and $a_0 \le 2^N$:
\[
  \forall k \le N,\quad \operatorname{orbit}(2^N - a_0, k) + \operatorname{antihydraOrbit}(a_0, k) = 3^k \cdot 2^{N - k},
\]
and for all $k < N$, $\operatorname{orbit}(2^N - a_0, k) \bmod 2 = \operatorname{antihydraOrbit}(a_0, k) \bmod 2$.
\end{theorem}

\begin{theorem}[Zero-Carry Base-3 Digit-Sum Conservation (SB.9)]
\emph{Lean identifier:} \nolinkurl{zero_carry_base3_digit_sum_and_even_step_invariance}.
Whenever $2^N g_N = 3^N g_0 + K_N$ with $K_N < 3^N$:
\[
  \operatorname{base3DigitSum}(2^N g_N) = \operatorname{base3DigitSum}(g_0) + \operatorname{base3DigitSum}(K_N),
\]
and whenever $\operatorname{orbit}(g_0, N) \bmod 2 = 0$, $\operatorname{base3DigitSum}(2^{N+1}\operatorname{orbit}(g_0, N+1)) = \operatorname{base3DigitSum}(2^N\operatorname{orbit}(g_0, N))$.
\end{theorem}

\begin{theorem}[Late-Step Cauchy Pinch \& Quotient GCD (SB.10)]
\emph{Lean identifier:} \nolinkurl{late_step_cauchy_pinch_and_odd_quotient_gcd_dvd}.
Across any late block segment $[N, N+M]$ of step sum $d = T_{N+M} - T_N$:
\[
  2^{d + s_{N+M}} u_{N+M} = 3^d \cdot 2^{s_N} u_N + A_{N, N+M},\qquad \gcd(u_N, u_{N+M}) \mid A_{N, N+M},
\]
and when $0 < \Sigma_M(s')$, $2(2^{d + s_{N+M}} u_{N+M} - 3^d \cdot 2^{s_N} u_N) < 3^d$.
\end{theorem}

\begin{theorem}[Affine $\times(3/2)$ Counter Conjugacy (SB.11)]
\emph{Lean identifier:} \nolinkurl{affine_counter_mahler_conjugacy_and_closed_bb6_non_halting}.
For every $m, c_0, k \in \mathbb{N}$, $\operatorname{affineCounterOrbit}(m, c_0, k) + 2m = \operatorname{orbit}(c_0 + 2m, k)$ and $\operatorname{affineCounterOrbit}(m, c_0, k) \bmod 2 = \operatorname{orbit}(c_0 + 2m, k) \bmod 2$.
\end{theorem}

\begin{theorem}[Dual Base-2/Base-3 Block Digit-Sum Laws (SB.12)]
\emph{Lean identifier:} \nolinkurl{dual_base2_base3_block_digitsum_conservation_and_parity}.
Across any \texttt{ValidBlockChain b\_0 s u (r + 2)} with $b_k = \operatorname{orbit}(b_0, T_k)$ and $P_k = 3 b_k + 1 = 2^{s_k} u_k$:
\[
  s_2(P_r) = s_2(u_r),\quad s_2(u_{r+1}) = s_2(3^{s_r} u_r + 1),\quad s_3(b_{r+1}) = s_3(u_r),\quad s_3(P_r) = s_3(b_r) + 1,\quad s_3(2^{s_{r+1}} u_{r+1}) = s_3(u_r) + 1,
\]
with $s_3(u_r) \bmod 2 = 1$, $s_3(P_r) \bmod 2 = 0$, and $s_3(P_r) \ne s_3(u_r)$.
\end{theorem}

\begin{theorem}[Weighted Base-3 Carry \& Senge--Straus Blueprint (SB.13--SB.14)]
\emph{Lean identifiers:} \nolinkurl{one_step_and_weighted_telescoping_base3_carry_identity}, \nolinkurl{top_down_senge_straus_stewart_joint_digitsum_blueprint}.
For every $g_0, k, N \in \mathbb{N}$:
\[
  s_3(2 g_{k+1}) = s_3(g_k) + (g_k \bmod 2),\qquad
  2^N s_3(g_N) = s_3(g_0) + \sum_{k=0}^{N-1} 2^k (g_k \bmod 2) + \sum_{k=0}^{N-1} 2^k\big(2 s_3(g_{k+1}) - s_3(2 g_{k+1})\big),
\]
and at every block $r$, $s_2(P_r) + s_3(P_{r+1}) = s_2(u_r) + s_3(u_r) + 1$.
\end{theorem}

\begin{theorem}[Legendre Factorial Syzygy \& Weight-2 Collapse (SB.15--SB.17)]
\emph{Lean identifiers:} \nolinkurl{legendre_factorial_valuation_dual_block_tower}, \nolinkurl{universal_mod4_quotient_rung_parity_lock_and_mersenne_trap}, \nolinkurl{mod1024_and_weight2_odd_quotient_three_block_domino_collapse}.
Across any \texttt{ValidBlockChain b\_0 s u (r + 2)}:
\begin{enumerate}
  \item $v_2(P_r!) - v_2(u_r!) = (2^{s_r} - 1) u_r$, $2(v_3(P_{r+1}!) - v_3(u_r!)) = (3^{s_r} - 1) u_r$, and $(3^{s_r} - 1)(v_2(P_r!) - v_2(u_r!)) = 2(2^{s_r} - 1)(v_3(P_{r+1}!) - v_3(u_r!))$.
  \item Whenever $2 \le s_{r+1}$, $(s_r \bmod 2 = 0 \iff u_r \bmod 4 = 3)$ and $(s_r \bmod 2 = 1 \iff u_r \bmod 4 = 1)$, with $s_r \bmod 2 = 0 \implies 2 \le s_2(u_r)$ and $(u_r \bmod 16 = 15 \wedge s_r = 2) \implies s_{r+1} = 3$.
  \item For every $u_K = 1024 w + 1$ (including all Hamming-weight-2 integers $u_K = 2^m + 1$ with $m \ge 10$), even $s_K$ dies in 1 block ($s_{K+1} = 1$) and $s_K = 3$ dies in 3 blocks ($s_{K+1} = 2, s_{K+2} = 6, s_{K+3} = 1$).
\end{enumerate}
\end{theorem}

\begin{theorem}[Mod-12 $\{5,7\}$ Lock \& Sharpened Ceiling (SB.18--SB.20)]
\emph{Lean identifiers:} \nolinkurl{universal_mod12_quotient_lock_and_joint_digitsum_ge_five}, \nolinkurl{one_step_backward_quotient_recovery_and_mod_three_pow_one}, \nolinkurl{sharpened_five_seven_archimedean_rung_ceiling}.
Across any \texttt{ValidBlockChain b\_0 s u (r + 2)} with $2 \le s_{r+1}$:
\begin{enumerate}
  \item $(s_r \bmod 2 = 0 \iff u_r \bmod 12 = 7)$ and $(s_r \bmod 2 = 1 \iff u_r \bmod 12 = 5)$; hence $5 \le u_r$ ($7 \le u_r$ when $s_r$ is even), $u_r \notin \{1, 3\}$, $2 \le s_2(u_r)$, $2 \le s_3(u_r)$, and $5 \le s_2(P_r) + s_3(P_{r+1})$.
  \item $\lfloor 2^{s_{r+1}} u_{r+1} / 3^{s_r} \rfloor = u_r$ and $(2^{s_{r+1}} u_{r+1}) \bmod 3^{s_r} = 1$.
  \item When $0 < \Sigma_r(s)$, $10 \cdot 2^{T_{r+1}} < 3^{T_r}(2 P_0 + 1)$ ($14 \cdot 2^{T_{r+1}} < 3^{T_r}(2 P_0 + 1)$ when $s_r$ is even).
\end{enumerate}
\end{theorem}

\begin{theorem}[Disjoint Base-3 Blocks \& Two-Sided CRT Pinch (SB.21--SB.23)]
\emph{Lean identifiers:} \nolinkurl{two_step_disjoint_base3_digit_block_law_and_quotient_recovery}, \nolinkurl{all_k_step_zero_carry_top_block_and_quotient_recovery}, \nolinkurl{two_sided_past_future_crt_pinch}.
Across any \texttt{ValidBlockChain b\_0 s u}:
\begin{enumerate}
  \item $2^{s_{r+1}+s_{r+2}} u_{r+2} = 3^{s_r+s_{r+1}} u_r + (3^{s_{r+1}} + 2^{s_{r+1}})$ with $\lfloor 2^{s_{r+1}+s_{r+2}} u_{r+2} / 3^{s_r+s_{r+1}} \rfloor = u_r$ and $s_3(2^{s_{r+1}+s_{r+2}} u_{r+2}) = s_3(u_r) + 1 + s_3(2^{s_{r+1}})$.
  \item Whenever $0 < \Sigma_M(s')$ with $s'(t) = s_{r+1+t}$, $d = \sum_{t=0}^{M-1} s_{r+1+t}$, $S = d + s_{r+1+M}$, $T = s_r + d$, and $R = 3^d + A_{r+1, r+1+M}$, we have $2^S u_{r+1+M} = 3^T u_r + R$ with $R < 2 \cdot 3^d < 3^T$, $\lfloor 2^S u_{r+1+M} / 3^T \rfloor = u_r$, and $s_3(2^S u_{r+1+M}) = s_3(u_r) + s_3(R)$.
  \item Every interior odd leftover $u_{r+1}$ satisfies $(2^{s_{r+1}} u_{r+1}) \bmod 3^{s_r} = 1$, $(3^{s_{r+1}} u_{r+1} + 1) \bmod 2^{s_{r+2}+1} = 2^{s_{r+2}}$, and $\gcd(3^{s_r}, 2^{s_{r+2}+1}) = 1$, with combined modulus $3^{s_r} \cdot 2^{s_{r+2}+1} \ge 72$ whenever $s_r, s_{r+2} \ge 2$.
\end{enumerate}
\end{theorem}

\begin{theorem}[Slack-Free $4/5$ Ceiling \& Sliding CRT Pinch (SB.24--SB.27)]
\emph{Lean identifiers:} \nolinkurl{unconditional_four_fifths_renewal_ceiling_and_past_base3_recovery}, \nolinkurl{unconditional_all_k_step_odd_leftover_recovery_and_nine_fifths_bound}, \nolinkurl{sliding_cut_point_past_future_crt_pinch}.
Across any \texttt{ValidBlockChain b\_0 s u}:
\begin{enumerate}
  \item Whenever $\forall k < r,\, 2 \le s_k$ (with no slack hypothesis needed), $5 A_{0,r} + 4 \cdot 2^{T_r} \le 4 \cdot 3^{T_r}$, $A_{0,r} < 3^{T_r}$, $\lfloor 2^{T_r} P_r / 3^{T_r} \rfloor = P_0$, $(2^{T_r} P_r) \bmod 3^{T_r} = A_{0,r}$, and $s_3(2^{T_r} P_r) = s_3(P_0) + s_3(A_{0,r})$.
  \item Whenever $\forall k < M,\, 2 \le s_{r+1+k}$, $R = 3^d + A_{r+1, r+1+M}$ satisfies $5 R + 4 \cdot 2^d \le 9 \cdot 3^d$, $R < 2 \cdot 3^d < 3^T$, $\lfloor 2^S u_{r+1+M} / 3^T \rfloor = u_r$, $(2^S u_{r+1+M}) \bmod 3^T = R$, and $s_3(2^S u_{r+1+M}) = s_3(u_r) + s_3(R)$.
  \item At every interior cut point $r$ with $\mathcal{M}_r = 3^{T_r} \cdot 2^{D_r + 1}$, $(2^{T_r} P_r) \bmod 3^{T_r} = A_{0,r}$, $(3^{d_r} P_r + A_{r, r+M}) \bmod 2^{D_r + 1} = 2^{D_r}$, $\gcd(3^{T_r}, 2^{D_r + 1}) = 1$, $5 \cdot 2^{T_r + D_r} P_r < 3^{T_r} \cdot 2^{D_r}(5 P_0 + 4)$, and whenever $P_0 < 2^{T_r + D_r + 1}$, $P_r < \mathcal{M}_r$.
  \item In addition, $A_{r+1, r+1+M} < 3^d$, $\lfloor R / 3^d \rfloor = 1$, $R \bmod 3^d = A_{r+1, r+1+M}$, $s_3(R) = 1 + s_3(A_{r+1, r+1+M})$, and $s_3(2^S u_{r+1+M}) = s_3(u_r) + 1 + s_3(A_{r+1, r+1+M})$ (\nolinkurl{zero_carry_base3_horner_remainder_ladder}).
\end{enumerate}
\end{theorem}

\begin{theorem}[Dual-Universe Series \& Base-3 Horner Evolution (SB.28--SB.30)]
\emph{Lean identifiers:} \nolinkurl{bridge_a_dual_universe_series_lowest_terms_and_uniqueness}, \nolinkurl{bridge_b_seed_independent_base3_step_evolution}, \nolinkurl{bridge_b_additive_carry_isolation_and_horner_step}.
Across any \texttt{ValidBlockChain b\_0 s u} with $s_k \ge 2$:
\begin{enumerate}
  \item $Q_K = (2^{T_{K+1}} u_K) / 3^{T_K} = P_0 + A_{0,K}/3^{T_K}$ is in strict lowest terms ($\gcd(2^{T_{K+1}} u_K, 3^{T_K}) = 1$), with $\lfloor 2^{T_{K+1}} u_K / 3^{T_K} \rfloor = P_0$ in $\mathbb{R}$, $(3^{T_K} P_0 + A_{0,K}) \bmod 2^{T_{K+1}+1} = 2^{T_{K+1}}$ in $\mathbb{Q}_2$, and unique modular root $P_0' = P_0$ below $2^{T_{K+1}}$ whenever $P_0 < 2^{T_{K+1}}$.
  \item For every $r \le K$, $s_3(2^{T_{r+1}} u_r) = s_3(P_0) + s_3(A_{0,r})$, so $s_3(2^{T_{r_2+1}} u_{r_2}) + s_3(A_{0,r_1}) = s_3(2^{T_{r_1+1}} u_{r_1}) + s_3(A_{0,r_2})$ is independent of $b_0$.
  \item For $X = 2^S u_{r+1+M} = 3^T u_r + R$, even after multiplying by $2^{s_r}$ we have $5(2^{s_r} R) < 4 \cdot 3^T$, $s_3(2^{s_r} X) = s_3(2^{s_r} u_r) + s_3(2^{s_r} R)$, and $s_3(2^{s_r} X) + s_3(u_r) + s_3(R) = s_3(X) + s_3(2^{s_r} u_r) + s_3(2^{s_r} R)$.
\end{enumerate}
\end{theorem}

\begin{theorem}[Binary Depletion \& Halving-Budget Conservation (SB.31--SB.33)]
\emph{Lean identifiers:} \nolinkurl{bridge_c_one_step_binary_trailing_ones_depletion_law}, \nolinkurl{bridge_c_multi_step_dual_base2_base3_zero_carry_split}, \nolinkurl{bridge_c_telescoped_binary_halving_budget_law}.
Across any \texttt{ValidBlockChain b\_0 s u}:
\begin{enumerate}
  \item $3^{s_r} u_r = 2^{s_{r+1}}(u_{r+1} - 1) + (2^{s_{r+1}} - 1)$ with $\lfloor 3^{s_r} u_r / 2^{s_{r+1}} \rfloor = u_{r+1} - 1$, $(3^{s_r} u_r) \bmod 2^{s_{r+1}} = 2^{s_{r+1}} - 1$, and $s_2(3^{s_r} u_r) + 1 = s_2(u_{r+1}) + s_{r+1}$.
  \item For $N_K = 3^{T_K} P_0 + A_{0,K} = 2^{T_{K+1}} u_K$ ($s_k \ge 2$), $N_K$ splits with zero carries in Base 3 ($s_3(N_K) = s_3(P_0) + s_3(A_{0,K})$) while $N_K - 1$ splits with zero carries in Base 2 ($s_2(N_K - 1) + 1 = T_{K+1} + s_2(u_K)$).
  \item Across $K$ blocks, $(T_{K+1} - s_0) + s_2(u_K) + \sum_{r=0}^{K-1} s_2(u_r) = s_2(u_0) + K + \sum_{r=0}^{K-1} s_2(3^{s_r} u_r)$.
\end{enumerate}
\end{theorem}

\begin{theorem}[Overnight 20-Loop Bridge Synthesis (SB.34--SB.53)]
\emph{Lean identifier:} \nolinkurl{grand_semigroup_breakout_53_bridge_synthesis}.
Across any \texttt{ValidBlockChain b\_0 s u (K + 4)} with $\forall k < K + 4,\, 2 \le s_k$, setting $N_K = 3^{T_K} P_0 + A_{0,K} = 2^{T_{K+1}} u_K$:
\begin{enumerate}
  \item \textbf{Kummer--Legendre Factorial Bridges (SB.34--SB.38):} $v_2((3^{s_r} u_r)!) + s_{r+1} = v_2(P_{r+1}!)$, $v_3(N_K!) = v_3((3^{T_K} P_0)!) + v_3(A_{0,K}!)$, $v_2(N_K!) + s_2(u_K) = 2 v_3(N_K!) + s_3(P_0) + s_3(A_{0,K})$, and $s_3(2^s m) + 2(v_3((2^s m)!) - 2^s v_3(m!)) = 2^s s_3(m)$.
  \item \textbf{Interior Mod-36 / Mod-24 / Mod-16 \& $79\times$ Ceiling Bridges (SB.39--SB.43):} $s_{r+1} \bmod 6 \in \{0,1,2,3,4,5\} \iff u_{r+1} \bmod 36 \in \{19,5,7,17,31,29\}$, $s_{r+1} \le 6 \implies 17 \le u_{r+1}$, $s_{r+1} = 2 \implies 79 \le u_{r+1}$, and $s_{r+1} = 2 \implies 395 \cdot 2^{T_{r+2}} < 3^{T_{r+1}}(5 P_0 + 4)$.
  \item \textbf{Wronskian, Cross-Ratio \& $4/5$ Hankel Sandwich Bridges (SB.44--SB.48):} $2^{s_r} 3^{s_{r+1}} P_{r+1}(P_{r+1}-1) = 3^{s_r} 2^{s_{r+1}} P_r(P_{r+2}-1)$ (so $s_r = s_{r+1} \implies P_{r+1}^2 - P_r P_{r+2} = P_{r+1} - P_r > 0$), $A_{0,r+M} = 3^d A_{0,r} + 2^{T_r}(2^d P_{r+M} - 3^d P_r)$, and $-4 \cdot 3^d P_3 < 5 \cdot 2^d(P_1 P_4 - P_2 P_3) < 4 \cdot 3^d P_1$.
  \item \textbf{Base-2 / Base-3 Complement Duality \& Grand Synthesis (SB.49--SB.53):} $s_2(2^N - 1 - m) + s_2(m) = N$ for $m < 2^N$ and $s_3(3^N - 1 - m) + s_3(m) = 2N$ for $m < 3^N$.
\end{enumerate}
\end{theorem}

\begin{theorem}[Collatz Syracuse Bridge \& Constant-Rung Bounds (SB.54--SB.56)]
\emph{Lean identifiers:} \nolinkurl{collatz_syracuse_bridge_and_one_two_cycle_exclusion}, \nolinkurl{universal_constant_rung_affine_invariance_and_run_length_bound}, \nolinkurl{sharp_four_fifths_equality_characterization_and_two_two_escape}.
\begin{enumerate}
  \item Every Collatz Syracuse step $2^s u' = 3 u + 1$ ($u'$ odd) satisfies $\lfloor 2^s u' / 3 \rfloor = u$, $(2^s u') \bmod 3 = 1$, $s_3(2^s u') = s_3(u) + 1$, and $s_2(3 u) + 1 = s_2(u') + s$; in addition, $s = 2, u_0 = 1$ is the unique Collatz 1-cycle and there are no Collatz 2-cycles with $1 \le u_0 < u_1$.
  \item Across $m$ consecutive Mahler blocks of constant rung $s \ge 1$, $2^{s m}((3^s - 2^s) u_{r+m} + 1) = 3^{s m}((3^s - 2^s) u_r + 1)$, forcing $2^{s m} \mid ((3^s - 2^s) u_r + 1)$ and $2^{s m} \le (3^s - 2^s) u_r + 1$.
  \item Equality $5 A_{0,r} + 4 \cdot 2^{T_r} = 4 \cdot 3^{T_r}$ holds whenever $\forall k < r,\, s_k = 2$, whereas $(s_0 \ge 2 \wedge s_1 \ge 3) \implies 5 A_{0,2} + 4 \cdot 2^{T_2} < 4 \cdot 3^{T_2}$, and at $s = (2, 2)$, $2 A_{0,2} = 104 > 81 = 3^4$ ($\Sigma_2(2,2) = -23 < 0$).
\end{enumerate}
\end{theorem}

\begin{theorem}[Diagonal Tail-Freeze \& $j \ge k$ Excess Laws (SB.57--SB.60)]
\emph{Lean identifiers:} \nolinkurl{two_window_jk_diagonal_cauchy_tail_freeze_and_quotient_invariance}, \nolinkurl{j_ge_k_excess_explosion_binary_swelling_and_halving_debt_law}, \nolinkurl{three_adic_seed_diagonal_shift_and_dead_zone_invariance_law}, \nolinkurl{grand_jk_diagonal_and_three_heavy_regime_60_synthesis}.
\begin{enumerate}
  \item \textbf{Two-Window Cauchy Tail-Freeze \& Quotient Invariance (SB.57):} Across consecutive windows of $M_1$ blocks ($n_1$ steps, ending rung $k_1$) and $M_2$ blocks ($n_2$ steps, ending rung $k_2$) starting from $j = s_r$, $2^{n_1+n_2+k_2} u_{r+1+M_1+M_2} = 3^{n_2}(2^{n_1+k_1} u_{r+1+M_1}) + 2^{n_1} A^{(2)} = 3^{j+n_1+n_2} u_r + R_{12}$ with $5(R_{12} - 3^{n_2} R_1) + 4 \cdot 2^{n_1+n_2} \le 4 \cdot 2^{n_1} \cdot 3^{n_2}$ and $\lfloor 2^{n_1+n_2+k_2} u_{r+1+M_1+M_2} / 3^{j+n_1+n_2} \rfloor = u_r = \lfloor 2^{n_1+k_1} u_{r+1+M_1} / 3^{j+n_1} \rfloor$.
  \item \textbf{$j \ge k$ Excess-Exponent Explosion \& Halving-Debt Crisis (SB.58):} Whenever $j = s_r \ge k = s_{r+1} \ge 2$ with $\Delta = j - k$, $2^k u_{r+1} = 3^\Delta \cdot 3^k u_r + 1$, $4 u_{r+1} \ge 9 \cdot 3^\Delta u_r + 1$, $u_{r+1} > 2 \cdot 3^\Delta u_r$, $s_2(u_{r+1}) + k = s_2(3^j u_r) + 1$, and across the next step $k' = s_{r+2}$, $2^{k+k'} u_{r+2} = 3^{j+k} u_r + 3^k + 2^k$, so $u_{r+2} \le u_r \implies 2^{k+k'} > 3^{j+k} \wedge 2^{k'} > 3^j$.
  \item \textbf{$3^{j_0}$ Diagonal Shift \& $j_0$-Trit Dead-Zone Law (SB.59):} Whenever $P_0 = 3^{j_0} m$ (or $u_r = 3^{j_0} m$), $N_K = 2^{n+k} u_K = 3^{j_0+n} m + A_{0,K}$ satisfies $\lfloor N_K / 3^{j_0+n} \rfloor = m$, $(N_K / 3^n) \bmod 3^{j_0} = 0$, $((N_K \bmod 3^{j_0+n}) / 3^n) = 0$, and $s_3(2^{n+k} u_K) = s_3(m) + s_3(A_{0,K})$.
\end{enumerate}
\end{theorem}

\begin{theorem}[Seed Finiteness \& Peak--Trough Seesaw (SB.61--SB.64)]
\emph{Lean identifiers:} \nolinkurl{local_descent_rung_jump_and_multi_step_descent_explosion}, \nolinkurl{trough_rebound_compound_interest_and_local_minimum_trap}, \nolinkurl{finite_seed_global_peak_trough_seesaw_and_return_barrier}, \nolinkurl{grand_seed_finiteness_and_extrema_64_synthesis}.
\begin{enumerate}
  \item \textbf{Local Descent Rung Jump \& Staircase Explosion (SB.61):} Whenever $u_{r+1} \le u_r$, we have $3^{s_r} < 2^{s_{r+1}}$ ($s_r \ge 2 \implies s_{r+1} \ge 4$, $s_r \ge 4 \implies s_{r+1} \ge 7$, $s_r \ge 7 \implies s_{r+1} \ge 12$), and across any $3$-step non-increasing run $u_{r+3} \le u_{r+2} \le u_{r+1} \le u_r$ with $s_r \ge 2$, $4 \le s_{r+1} \wedge 7 \le s_{r+2} \wedge 12 \le s_{r+3}$.
  \item \textbf{Trough-Rebound Compound-Interest Trap (SB.62):} Whenever $2 \le s_r$ and $u_{r+1} \le u_r$, $4 \le s_{r+1}$, $81 u_{r+1} + 1 \le 2^{s_{r+2}} u_{r+2}$, and $81(3^{s_r} u_r + 1) + 16 \le 16 \cdot 2^{s_{r+2}} u_{r+2}$; if also $2 \le s_{r+2} \le s_{r+1}$, then $9 \cdot 3^{s_{r+1}-s_{r+2}} u_{r+1} + 1 \le 4 u_{r+2}$.
  \item \textbf{Finite-Seed Global Peak--Trough Seesaw \& Post-Horizon Uniqueness (SB.63):} Across $K \ge 1$ blocks with $\forall k < K,\, 2 \le s_k$, $3^{T_K} \cdot 2^{s_0} u_0 < 2^{T_K + s_K} u_K$ and $5 \cdot 2^{T_K + s_K} u_K + 4 \cdot 2^{T_K} \le 3^{T_K}(5 \cdot 2^{s_0} u_0 + 4)$; consequently $u_K \le u_0 \implies 3^{T_K} \cdot 2^{s_0} < 2^{T_K + s_K}$ and $s_K \le s_0 \implies 3^{T_K} u_0 < 2^{T_K} u_K$.
\end{enumerate}
\end{theorem}

\begin{theorem}[$\text{LTE}_2$ Flush Pinning \& Baker--Matveev Pinch (SB.65--SB.68)]
\emph{Lean identifiers:} \nolinkurl{lte2_ultrametric_flush_pinning_and_descent_prohibition}, \nolinkurl{post_horizon_binary_zero_bit_desert_and_dual_radix_lock}, \nolinkurl{shifted_basepoint_baker_matveev_pinch_and_repeated_core_rhin_form}, \nolinkurl{senge_straus_stewart_finiteness_bridge_and_grand_68_synthesis}.
\begin{enumerate}
  \item \textbf{$\text{LTE}_2$ Ultrametric Flush Pinning (SB.65):} For $s_r, s_{r+1} \ge 2$, $(s_r \bmod 4 = 2) \implies (u_r \bmod 16 = 15 \leftrightarrow s_{r+1} = 3) \wedge (u_r \bmod 16 = 7 \leftrightarrow 4 \le s_{r+1}) \wedge (u_{r+1} \le u_r \to u_r \bmod 16 = 7)$, with exact companion characterizations for $s_r \bmod 4 \in \{0, 1, 3\}$.
  \item \textbf{Post-Horizon Binary Zero-Bit Desert (SB.66):} Whenever $P_0 < 2^{B_0}$ and $B_0 \le S_K = T_K + s_K$, $\forall B \ge B_0,\, \lfloor (P_0 \bmod 2^{S_K}) / 2^B \rfloor = 0$, $s_2(N_K) = s_2(u_K)$, and $s_3(N_K) = s_3(P_0) + s_3(A_{0,K})$.
  \item \textbf{Shifted-Basepoint Baker--Matveev Pinch \& Repeated-Core Rhin Reduction (SB.67):} Across $[r, r+M]$ of step sum $n_2$, $5(2^{n_2+s_K} u_{r+M} - 3^{n_2} 2^{s_r} u_r) + 4 \cdot 2^{n_2} \le 4 \cdot 3^{n_2}$ with $3^{T_r} 2^{s_0} u_0 < 2^{T_r}(2^{s_r} u_r)$; whenever $u_{r+M} = u_r$, $(2^{s_r} u_r)(2^{n_2+s_K-s_r} - 3^{n_2}) = A_{r,M}$, forcing $5(2^{s_r} u_r) + 4 \cdot 2^{n_2} \le 4 \cdot 3^{n_2}$ and $5 \cdot 3^{T_r} 2^{s_0} u_0 < 4 \cdot 2^{T_r} 3^{n_2}$.
  \item \textbf{Senge--Straus / Stewart Finiteness Bridge (SB.68):} $s_2(N_K) = s_2(u_K)$, $s_3(N_K) = s_3(P_0) + s_3(A_{0,K})$, and $s_3(N_K) + 2(v_3(N_K!) - 2^{S_K} v_3(u_K!)) = 2^{S_K} s_3(u_K)$.
\end{enumerate}
\end{theorem}

\begin{theorem}[Upgraded Baker Pinch \& $1/5$ Reduced Ceiling (SB.69--SB.72)]
\emph{Lean identifiers:} \nolinkurl{adjacent_odd_core_coprimality_and_universal_non_equality}, \nolinkurl{two_step_odd_core_gcd_and_return_extinction}, \nolinkurl{end_of_block_three_pow_upgraded_baker_pinch_and_s_r_free_return_bound}, \nolinkurl{exact_two_adic_renewal_valuation_and_one_fifth_reduced_ceiling}.
\begin{enumerate}
  \item \textbf{Adjacent Coprimality \& Universal Non-Equality (SB.69):} Whenever $2 \le s_r$, $\gcd(u_r, u_{r+1}) = 1$, $\gcd(2^{s_{r+1}} u_{r+1}, 3^{s_r} u_r) = 1$, and $u_{r+1} \ne u_r$.
  \item \textbf{Two-Step Odd-Core GCD \& $s_r$-Free Return Bound (SB.70):} $\gcd(u_r, u_{r+2}) \mid (3^{s_{r+1}} + 2^{s_{r+1}})$, and whenever $u_{r+2} = u_r$, $(2^{s_{r+1}+s_{r+2}} - 3^{s_r+s_{r+1}}) u_r = 3^{s_{r+1}} + 2^{s_{r+1}}$, $u_r \mid (3^{s_{r+1}} + 2^{s_{r+1}})$, and $u_r \le 3^{s_{r+1}} + 2^{s_{r+1}}$.
  \item \textbf{End-of-Block $3^{s_r}$-Upgraded Baker Pinch (SB.71):} For $R = 3^d + A_{r+1, M}$, $R \bmod 2 = 1$, $\gcd(u_r, u_{r+1+M}) \mid R$, $5 \cdot 3^{s_r}(2^S u_{r+1+M} - 3^T u_r) + 4 \cdot 3^{s_r} 2^d \le 9 \cdot 3^T$, and $u_{r+1+M} = u_r \implies 5 u_r + 4 \cdot 2^d \le 9 \cdot 3^d$.
  \item \textbf{Exact 2-Adic Renewal Valuation \& $1/5$ Reduced Ceiling (SB.72):} $A_{0,K} = 2^{s_0} R_0$ with $R_0 \bmod 2 = 1$, $2^{S_0} u_K = 3^{T_K} u_0 + R_0$, and $5 R_0 + 4 \cdot 2^{d_0} \le 3^{T_K}$ ($5 R_0 < 3^{T_K}$) for all $s_0 \ge 2$.
\end{enumerate}
\end{theorem}

\begin{theorem}[Core Subordinacy, Product Formula \& Ridout Bridge (SB.73--SB.76)]
\emph{Lean identifiers:} \nolinkurl{cubic_and_quintic_odd_core_subordinacy_to_two_pow}, \nolinkurl{global_multiplicative_to_additive_product_formula_bridge}, \nolinkurl{two_window_ridout_subspace_tail_cancellation_bridge}, \nolinkurl{grand_global_dominance_and_product_formula_76_synthesis}.
\begin{enumerate}
  \item \textbf{Cubic $(3,5)$ \& Quintic $(5,8)$ Odd-Core Subordinacy (SB.73):} Across $K = M + 1 \ge 1$ blocks with $s_k \ge 2$, $125 \cdot 2^{3 s_K} u_K^3 < 2^{2 T_K + 3 s_0}(5 u_0 + 1)^3$ and $3125 \cdot 2^{5 s_K} u_K^5 < 2^{3 T_K + 5 s_0}(5 u_0 + 1)^5$.
  \item \textbf{Global Multiplicative-to-Additive Product Formula (SB.74):} $2^{T_K+s_K} u_K \prod_{r=0}^{K-1} u_r = (3^{T_K} P_0 + A_{0,K}) \prod_{r=0}^{K-1} u_r = P_0 \prod_{r=0}^{K-1} (3^{s_r} u_r + 1)$, forcing $P_0 \mid (A_{0,K} \prod_{r=0}^{K-1} u_r)$.
  \item \textbf{Two-Window Ridout/Subspace Cancellation Bridge (SB.75):} For $\Delta = 2^{T_r+n_2+s_K} u_{r+M} - 3^{n_2}(2^{T_r+s_r} u_r) = 2^{T_r} A_{r,M}$, $5 \Delta + 4 \cdot 2^{T_r+n_2} \le 4 \cdot 2^{T_r} 3^{n_2}$, and whenever $C_1 2^{T_r+s_r} 3^{n_2} \le C_2 3^E \Delta$ ($C_2 \ge 1$), both $2^{T_r}$ and $3^{n_2}$ cancel to yield $5 C_1 2^{s_r} < 4 C_2 3^E$.
\end{enumerate}
\end{theorem}

\begin{theorem}[Greedy Base-$(2/3)$ Kempner Tail Valuation (SB.77--SB.80)]
\emph{Lean identifiers:} \nolinkurl{greedy_two_adic_base_two_thirds_kempner_step_and_prefix_law}, \nolinkurl{two_horizon_kempner_tail_exact_valuation_and_one_fifth_bound}, \nolinkurl{kempner_rational_target_obstruction_and_automaton_irrationality_bridge}, \nolinkurl{grand_kempner_two_adic_expansion_80_synthesis}.
\begin{enumerate}
  \item \textbf{Greedy 2-Adic Base-$(2/3)$ Step \& Prefix Law (SB.77):} For every $r \le K$, $3^{T_r} P_0 + A_{0,r} = 2^{T_{r+1}} u_r$ and $(3^{T_r} P_0 + A_{0,r}) \bmod 2^{T_{r+1}+1} = 2^{T_{r+1}}$; for every $r < K$, $3^{s_r}(3^{T_r} P_0 + A_{0,r}) + 2^{T_{r+1}} = 3^{T_{r+1}} P_0 + A_{0,r+1} = 2^{T_{r+2}} u_{r+1}$.
  \item \textbf{Two-Horizon Kempner Tail Valuation \& $1/5$ Ceiling (SB.78):} Between $K_1 = r$ and $K_2 = r + 1 + M$, $2^{S_2} u_{K_2} - 3^{s_r+d}(2^{T_r+s_r} u_{K_1}) = 2^{T_r+s_r} R_{\text{tail}}$ with $R_{\text{tail}} \bmod 2 = 1$ and $5 R_{\text{tail}} + 4 \cdot 2^d \le 3^{s_r+d}$.
  \item \textbf{Kempner Rational-Target Obstruction Bridge (SB.79):} $\gcd(2^{S_K} u_K, 3^{T_K}) = 1$, $\lfloor 2^{S_K} u_K / 3^{T_K} \rfloor = P_0$, $(3^{T_K} P_0 + A_{0,K}) \bmod 2^{S_K+1} = 2^{S_K}$, and $\forall M_0,\, (3^{T_K} M_0 + A_{0,K}) \bmod 2^{S_K} \ne 0 \implies P_0 \ne M_0$.
\end{enumerate}
\end{theorem}

\begin{theorem}[The $(5P+4)/36$ Ceiling-Dual Gauge (SB.81--SB.84)]
\emph{Lean identifiers:} \nolinkurl{ceiling_dual_gauge_transform_and_rung_two_parity_detector}, \nolinkurl{multi_block_gauge_deficit_master_identity_and_one_ninth_bound}, \nolinkurl{gauge_complement_base3_conservation_and_top_22_trit_lock}, \nolinkurl{dual_wall_monotone_squeeze_and_grand_84_synthesis}.
\begin{enumerate}
  \item \textbf{Ceiling-Dual Gauge \& Rung-2 Parity Detector (SB.81):} For every interior block $r + 1$ ($s_r, s_{r+1}, s_{r+2} \ge 2$), $P_{r+1} \bmod 36 = 28$, $36 W_{r+1} = 5 P_{r+1} + 4$, $(W_{r+1} \bmod 2 = 0 \leftrightarrow s_{r+1} = 2)$, $(W_{r+1} \bmod 2 = 1 \leftrightarrow 3 \le s_{r+1})$, and $2^{s_{r+1}} W_{r+2} + (3^{s_{r+1}-2} - 2^{s_{r+1}-2}) = 3^{s_{r+1}} W_{r+1}$ (so $s_{r+1} = 2 \implies 4 W_{r+2} = 9 W_{r+1}$).
  \item \textbf{Multi-Block Gauge Deficit \& $1/9$ Bound (SB.82):} Across $M \ge 1$ blocks of step sum $d$, $5 A_M + 4 \cdot 2^d + 36 E_M = 4 \cdot 3^d$, $2^d W_{r+1+M} + E_M = 3^d W_{r+1}$, $9 E_M + 2^d \le 3^d$, and $E_M < 3^{d-2}$.
  \item \textbf{Gauge-Complement Base-3 Conservation \& \texttt{22}-Trit Lock (SB.83--SB.84):} Whenever $1 \le E_M$, $\lfloor 2^d W_{r+1+M} / 3^d \rfloor = W_{r+1} - 1$, $(2^d W_{r+1+M}) \bmod 3^d = 3^d - E_M$, $\lfloor ((2^d W_{r+1+M}) \bmod 3^d) / 3^{d-2} \rfloor = 8$, and $s_3(2^d W_{r+1+M}) + s_3(E_M - 1) = s_3(W_{r+1} - 1) + 2 d$.
\end{enumerate}
\end{theorem}

\begin{theorem}[Super-Block Renormalization \& Mod-8 Predictor (SB.85--SB.88)]
\emph{Lean identifiers:} \nolinkurl{gauge_odd_valuation_trap_and_mod_four_survival_lock}, \nolinkurl{mod_eight_odd_gauge_future_rung_predictor_and_seven_extinction}, \nolinkurl{universal_cyclotomic_gauge_spectrum_and_two_three_eliminator}, \nolinkurl{grand_superblock_and_multigauge_88_synthesis}.
\begin{enumerate}
  \item \textbf{Odd-Valuation Death Trap \& $s=2$ Run Evolution (SB.85):} $W_{r+1} \bmod 4 \ne 2$, $(s_{r+1} = 2 \leftrightarrow W_{r+1} \bmod 4 = 0)$, $4^m W_{r+1+m} = 9^m W_{r+1}$ across $m$ blocks of $s=2$, and $W_{r+1} \ne 2 \cdot 4^m \cdot V$ for every odd $V$.
  \item \textbf{Mod-8 Future-Rung Predictor \& Super-Block Jump (SB.86):} For $W_{r+1} = 4^m V$ followed by $s^* = s_{r+1+m} \ge 3$, $W_{r+1+m} = 9^m V$, $V \bmod 2 = 1$, $V \bmod 8 \ne 7$, $(V \bmod 8 = 3 \leftrightarrow s^* = 3)$, $(V \bmod 8 = 5 \leftrightarrow s^* = 4)$, $(V \bmod 8 = 1 \leftrightarrow s^* \ge 5)$, and $2^{s^*} W_{r+1+m+1} + (3^{s^*-2} - 2^{s^*-2}) = 3^{2m+s^*} V$.
  \item \textbf{Universal Cyclotomic Gauge Spectrum (SB.87):} For $\mathcal{G}_r^{(m)} = (3^m - 2^m) P_r + 2^m$, $s_r = m \implies 2^m \mathcal{G}_{r+1}^{(m)} = 3^m \mathcal{G}_r^{(m)}$, and $19 \mathcal{G}_r^{(2)} = 5 \mathcal{G}_r^{(3)} + 36$.
\end{enumerate}
\end{theorem}

\begin{theorem}[Cyclotomic Velocity \& Canyon Interval Tube (SB.89--SB.92)]
\emph{Lean identifiers:} \nolinkurl{adaptive_cyclotomic_velocity_factorization_and_eliminator}, \nolinkurl{velocity_two_adic_valuation_lock_and_telescoped_odd_core_sum}, \nolinkurl{nested_canyon_interval_tube_and_exponential_width_contraction}, \nolinkurl{master_conditional_canyon_closure_and_92_theorem_synthesis}.
\begin{enumerate}
  \item \textbf{Adaptive Cyclotomic Velocity Factorization (SB.89):} For $K_r = (3^{s_r} - 2^{s_r}) u_r + 1 = P_{r+1} - P_r$, $\mathcal{G}^{(s_r)}(P_r) = 2^{s_r} K_r$, $\mathcal{G}^{(s_r)}(P_{r+1}) = 3^{s_r} K_r$, and $s_{r+1} = s_r \implies 2^{s_r} K_{r+1} = 3^{s_r} K_r$.
  \item \textbf{Velocity 2-Adic Valuation Lock \& Telescoped Sum (SB.90):} Whenever $s_r < s_{r+1}$, $K_r \bmod 2^{s_r+1} = 2^{s_r}$ ($v_2(K_r) = s_r$); whenever $s_{r+1} < s_r$, $K_r \bmod 2^{s_{r+1}+1} = 2^{s_{r+1}}$ ($v_2(K_r) = s_{r+1}$); and $P_K = P_0 + \sum_{r=0}^{K-1} K_r$.
  \item \textbf{Nested Canyon Interval Tube \& Master Conditional Closure (SB.91--SB.92):} For $L_{\text{num}}(k) = 5 \cdot 2^{T_k} P_k$ and $U_{\text{num}}(k) = 2^{T_k}(5 P_k + 4)$, $U_{\text{num}}(k) - L_{\text{num}}(k) = 4 \cdot 2^{T_k}$ and $3^d L_{\text{num}}(r) < L_{\text{num}}(r+M) < U_{\text{num}}(r+M) \le 3^d U_{\text{num}}(r)$, with exact conditional termination under either the Carry-Mixing Sparsity Obstruction or the Product-Formula Divisibility Obstruction.
\end{enumerate}
\end{theorem}

\begin{theorem}[Binary Carry-Split \& Reduced Product Formula (SB.93--SB.96)]
\emph{Lean identifiers:} \nolinkurl{path_a_binary_carry_split_and_quotient_decomposition_of_u_K}, \nolinkurl{path_a_zero_carry_binary_block_split_and_weight_excess_law}, \nolinkurl{path_b_reduced_odd_seed_product_formula_and_symmetric_excess}, \nolinkurl{grand_path_a_and_path_b_96_theorem_synthesis}.
\begin{enumerate}
  \item \textbf{Path A Binary Carry-Split \& Quotient Decomposition (SB.93):} In $2^{S_0} u_K = 3^{T_K} u_0 + R_0$, $(3^{T_K} u_0) \bmod 2^{S_0} + R_0 \bmod 2^{S_0} = 2^{S_0}$, $u_K = \lfloor 3^{T_K} u_0 / 2^{S_0} \rfloor + \lfloor R_0 / 2^{S_0} \rfloor + 1$, and $s_2((3^{T_K} u_0) \bmod 2^{S_0}) + s_2(R_0 \bmod 2^{S_0}) = S_0 + 1$.
  \item \textbf{Path A Full Binary Weight Excess Law (SB.94):} $s_2(3^{T_K} u_0) + s_2(R_0) = S_0 + 1 + s_2(\lfloor 3^{T_K} u_0 / 2^{S_0} \rfloor) + s_2(\lfloor R_0 / 2^{S_0} \rfloor)$, forcing $s_2(3^{T_K} u_0) + s_2(R_0) \ge S_0 + 2$ whenever $u_K \ge 2$.
  \item \textbf{Path B Reduced Odd-Seed Product Formula (SB.95):} Canceling $P_0 = 2^{s_0} u_0$ yields $(3^{T_K} u_0 + R_0) \prod_{r=1}^{K-1} u_r = \prod_{r=0}^{K-1} (3^{s_r} u_r + 1)$, $R_0 \prod_{r=1}^{K-1} u_r = \prod_{r=0}^{K-1} (3^{s_r} u_r + 1) - 3^{T_K} \prod_{r=0}^{K-1} u_r$, and $\prod_{r=0}^{K-1} (3^{s_r} u_r + 1) \bmod u_0 = (R_0 \prod_{r=1}^{K-1} u_r) \bmod u_0$.
\end{enumerate}
\end{theorem}

\begin{theorem}[Mod-40 Six-Residue Bijection \& Base-$D_m$ Locks (SB.97--SB.100)]
\emph{Lean identifiers:} \nolinkurl{universal_mod5_and_mod10_gauge_lock_and_run_parity_detector}, \nolinkurl{mod40_six_residue_superblock_bijection_and_34_class_extinction}, \nolinkurl{cyclotomic_base_dm_modular_invariance_and_mod19_mod13_locks}, \nolinkurl{grand_centennial_100_theorem_semigroup_breakout_synthesis}.
\begin{enumerate}
  \item \textbf{Universal Mod-5 \& Mod-10 Gauge Lock (SB.97):} At every interior block $r + 1$ ($s_r, s_{r+1}, s_{r+2} \ge 2$), $W_{r+1} \bmod 5 = 4$, $(W_{r+1} \bmod 10 = 4 \leftrightarrow s_{r+1} = 2)$, and $(W_{r+1} \bmod 10 = 9 \leftrightarrow 3 \le s_{r+1})$; for $W_{r+1} = 4^m V$ ending at $s^* \ge 3$, $(m \bmod 2 = 0 \leftrightarrow V \bmod 5 = 4 \leftrightarrow V \bmod 10 = 9)$ and $(m \bmod 2 = 1 \leftrightarrow V \bmod 5 = 1 \leftrightarrow V \bmod 10 = 1)$.
  \item \textbf{Mod-40 Six-Residue Super-Block Bijection (SB.98):} Across any Super-Block $W_{r+1} = 4^m V$ ending at $s^* \ge 3$, $V \bmod 40 \in \{1, 9, 11, 19, 21, 29\}$ (34 of 40 residue classes are extinct), with $V \bmod 40 \in \{19, 11, 29, 21, 9, 1\}$ in exact 6-way bijection with $(m \bmod 2) \times (s^* \in \{3, 4, \ge 5\})$.
  \item \textbf{Cyclotomic Base-$D_m$ \& Base-4/Base-9 Shift Laws (SB.99--SB.100):} $\mathcal{G}^{(m)}(P) \bmod (3^m - 2^m) = 2^m \bmod (3^m - 2^m)$ ($\mathcal{G}^{(2)}(P) \bmod 5 = 4$, $\mathcal{G}^{(3)}(P) \bmod 19 = 8$, $\mathcal{G}^{(4)}(P) \bmod 13 = 3$, $\mathcal{G}^{(5)}(P) \bmod 211 = 32$) and $s_3(9^m V) = s_3(V)$.
\end{enumerate}
\end{theorem}

\begin{theorem}[Fast Super-Block Sieve \& Pre/Post-Horizon Bridge (SB.101--SB.104)]
\emph{Lean identifiers:} \nolinkurl{odd_seed_block_zero_and_block_one_exact_extinction_classes}, \nolinkurl{mod40_window_thirty_four_of_forty_extinction_and_dyadic_uniqueness}, \nolinkurl{fast_superblock_sieve_macro_step_and_record_survivor_certificates}, \nolinkurl{zfc_deductive_semigroup_pre_vs_post_horizon_overspill_bridge}.
\begin{enumerate}
  \item \textbf{Block-0 \& Block-1 Unconditional Extinction Classes (SB.101):} For odd $b_0$, $b_0 \bmod 4 = 3 \implies s_0 = 1$ ($50\%$ die at Block 0), and $b_0 \bmod 16 = 1 \vee b_0 \bmod 32 = 29 \vee b_0 \bmod 64 = 5 \implies \neg(2 \le s_1)$ (die at Block 1).
  \item \textbf{Mod-40 Window $34/40$ Extinction \& Dyadic Uniqueness (SB.102):} Out of $0\dots 39$, exactly $6$ residues satisfy \texttt{isMod40Survivor} and $34$ are extinct; for $P_0 < 2^{S_K}$, at most one $P_0' < 2^{S_K}$ satisfies $(3^{T_K} P_0' + A_K) \bmod 2^{S_K} = 0$.
  \item \textbf{Fast $O(1)$ Super-Block Sieve Jump \& Record Certificates (SB.103):} $W_{r+1} = 4^m V \implies P_{r+1+m} = (36 \cdot 9^m V - 4)/5$, certifying the multi-block jumps of $b_0 = 13$ ($21$ steps) and $b_0 = 485$ ($25$ steps, skipping five $s=2$ blocks via $4^5 \to 9^5$).
  \item \textbf{ZFC Deductive Semigroup Pre- vs. Post-Horizon Bridge (SB.104):} For every $q \in \mathbb{N}$, $(3^{T_K}(P_0 + q 2^{S_K+1}) + A_K) \bmod 2^{S_K+1} = 2^{S_K}$ (pre-horizon periodic lattice), whereas $P_0 < 2^{S_K}$ excludes all $q \ge 1$ (post-horizon archimedean collapse).
\end{enumerate}
\end{theorem}

\section{Generalized $q/2$ Phase Diagram, Parry Cylinder Spectral Bounds, and Real-Companion Bridges}

\begin{theorem}[Generalized $q/2$ Kempner Inversion \& Core Gauge (SB.105--SB.108)]
\emph{Lean identifiers:} \nolinkurl{universal_q_half_zone_thresholds_and_kempner_inversion_dichotomy}, \nolinkurl{universal_q_half_mod4_one_step_extinction_and_four_adic_bijection}, \nolinkurl{universal_q_half_ceiling_dual_gauge_and_superblock_law}, \nolinkurl{grand_universal_q_half_relaxation_and_mahler_uniqueness_108_synthesis}.
\begin{enumerate}
  \item \textbf{Kempner Threshold Inversion (SB.105):} For every $q \ge 3$, $1/q < (q+2)/q^2 < 1/(q-2)$ and $q/(q^2-4) < 1/(q-2)$. At $q = 3$, $(3+2)(3^2-4) = 25 < 27 = 3^3$ ($5/9 < 3/5$, Kempner overflow $8/5 > 1$ at $k = 1/2$), whereas for every $q = d + 4 \ge 4$, $(q+2)(q^2-4) = q^3 + 2(d^2 + 6 d + 4) \ge q^3 + 8 > q^3$ ($q/(q^2-4) < (q+2)/q^2$).
  \item \textbf{One-Step Mod-4 50/50 Split \& Mod-$2^{S_0}$ Uniqueness (SB.106):} For odd $q = 2m+1$ and $Q = q^{s_r}$, among every consecutive odd pair $\{4t+1, 4t+3\}$ exactly one satisfies $(Q u + 1) \bmod 4 = 0$ ($s_{r+1} \ge 2$) and one satisfies $(Q u + 1) \bmod 4 = 2$ ($s_{r+1} = 1$), with unique seed recovery $P_0 < 2^{S_0}$ modulo $2^{S_0}$.
  \item \textbf{Core Gauge $(q^2 - 4) u + 1$ \& Exit Predictor (SB.107--SB.108):} For $q = 2k+1 \ge 3$, $4 u_{r+1} = q^2 u_r + 1 \implies 4((q^2-4)u_{r+1}+1) = q^2((q^2-4)u_r+1)$. Odd 2-adic valuation $(q^2-4)u_m + 1 = 2 q^{2m} V$ ($V$ odd) forces $(q^2 u_m + 1) \bmod 4 = 2$ ($s_{r+m+1} = 1$), while $(q^2-4) Y_m + 1 = q^{2m+2} V$ ($4 Y_m = q^2 u_m + 1$) locks $s^* \in \{3, 4, \ge 5\}$ from $V \bmod 8 \in \{3, 7\}, 5, 1$.
\end{enumerate}
\end{theorem}

\begin{theorem}[Adelic Odd-Modulus Inverse, Parry Tree \& Base-6 Law (SB.109--SB.114)]
\emph{Lean identifiers:} \nolinkurl{adelic_profinite_odd_modulus_invertibility_and_archimedean_prime_floor}, \nolinkurl{real_first_parry_transfer_map_invariance_and_dual_series_obstruction}, \nolinkurl{unconditional_effective_fibonacci_survivor_power_saving_bound}, \nolinkurl{parry_cylinder_tree_three_pow_sum_and_flp_effective_survivor_bound}, \nolinkurl{post_horizon_top_bit_desert_and_two_word_difference_valuation}, \nolinkurl{base_six_terminating_digit_characterization_and_kempner_depth_law}.
\begin{enumerate}
  \item \textbf{Odd-Modulus Solvability \& Divisor Floor (SB.109):} $(2^T (m+1)^T) \bmod (2m+1) = 1$ yields a unique root $P_{\text{sol}} \in [0, 2m+1)$ for every $P_0, A_K$, whereas in $\mathbb{N}$, $d \mid u_K$ forces $d \le u_K$ and $d \cdot 2^{s_K} \le P_K$.
  \item \textbf{Parry Transfer Invariance \& Dual Series (SB.110):} For $1 \le D$ and $2 R_0 < D$, $2 R_N < 2^N D$ and $R_N + D \Psi_N = 3^N R_0$; matching parities give $2^N g_N = 3^N b_0 + \Psi_N$; and $(3^L - 2^L) b_0 + \Psi_L = 0 \implies b_0 = 0 \wedge \Psi_L = 0$ ($L \ge 1$).
  \item \textbf{Fibonacci \& Parry $(3/2)^W$ Bounds (SB.111--SB.113):} On $[0, 2^W)$, \texttt{11}-free seeds satisfy $|\mathcal{S}^{\text{Fib}}_W| \le F_{W+2}$ ($F_{4k+2} \le 2^{3k+1}$), while \nolinkurl{parryCylTree} satisfies $\sum_{w \in \mathcal{C}_W} B_w = 3^W$ and $2^W |\mathcal{S}_W| + 2 \cdot 2^W \le 3^{W+1}$ for rational-companion survivor sets~\cite{Flatto1992,FLP1995}, with $n_0 = v_2(b_2 - b_1) < W$ for $b_1 < b_2 < 2^W$.
  \item \textbf{Base-6 Joint Characterization (SB.114):} For $g = \lfloor X/D \rfloor$, $R = X \bmod D$, $d_0 = g \bmod 6$, and $f_1 = \lfloor 6 R / D \rfloor$, $2 R < D \leftrightarrow f_1 < 3$, and $3(2 \cdot 4^T) 3^T = 6^{T+1} 2^T$.
\end{enumerate}
\end{theorem}

\begin{theorem}[Catalan--Pillai, Hankel Vanishing \& Word-Powers (SB.115--SB.120)]
\emph{Lean identifiers:} \nolinkurl{two_superblock_catalan_pillai_transition_and_coprime_odd_cores}, \nolinkurl{superblock_abc_radical_compression_and_conditional_run_cutoff}, \nolinkurl{higher_order_hankel_row_reduction_and_superblock_vanishing_laws}, \nolinkurl{universal_local_word_power_gauge_divisibility_and_cutoff}, \nolinkurl{rational_remainder_recurrence_integer_escape_and_shadowing_ball_cutoff}, \nolinkurl{grand_all_eight_angles_120_theorem_semigroup_breakout_synthesis}.
\begin{enumerate}
  \item \textbf{Two-Superblock Catalan--Pillai \& Prime Compression (SB.115--SB.116):} Consecutive Super-Blocks satisfy $2^{2m_2+s^*} V_2 + (3^{s^*-2} - 2^{s^*-2}) = 3^{2m_1+s^*} V_1$ and $\gcd(V_1, V_2) \mid (3^{s^*-2} - 2^{s^*-2})$, giving $3^{2m_1+3} V_1 = 2^{2m_2+3} V_2 + 1$ at $s^* = 3$, $\gcd(V_1, V_2) = 1$ at $s^* \in \{3, 4\}$, and prime support in $\{2, 3\} \cup \{p \mid D^* V_1 V_2\}$.
  \item \textbf{Higher-Order Hankel Vanishing (SB.117):} Whenever $2 g_{k+1} = 3 g_k + e_k$, $H_n^{(3)} = 0$ on $(0,0,0,0)$ and $H_n^{(4)} = 0$ identically whenever $(e_n, \dots, e_{n+5})$ has period $2$.
  \item \textbf{Local Word-Power $V^m$ Cutoff (SB.118):} Across $m$ repeats of a word $V$ of step sum $d \ge 1$ ($5 A_V \le 4(3^d - 2^d)$) along $\forall k < r,\, 2 \le s_k$, $2^{m d} \mid ((3^d - 2^d) P_r + A_V)$ and $5 \cdot 2^{T_r + m d} < 3^{T_r + d}(5 P_0 + 4)$ (retaining the seed factor $5 P_0 + 4$).
  \item \textbf{Rational Shadowing Escape (SB.119--SB.120):} $(3^{a+1}(2^a(2t+1))) \bmod 2^{a+1} = 2^a$, and whenever $X_n = (3^n p) \bmod (2^n q) \in [2^{n-1} q + E,\, 2^n q - E]$, every $A/B$ with $|q A - p B| 3^n < E B$ satisfies $2^{n-1} B < (3^n A) \bmod (2^n B)$.
\end{enumerate}
\end{theorem}

\begin{theorem}[Parry Entropy Sandwich \& Parseval Floor (SB.121--SB.124)]
\emph{Lean identifiers:} \nolinkurl{two_sided_parry_entropy_sandwich_and_nonsplit_lower_bound}, \nolinkurl{splitting_pair_half_period_antiphase_and_odd_parseval_identity}, \nolinkurl{conditional_fourier_sub_flp_exponent_and_seventeen_twelfths_crossover}, \nolinkurl{barrier_six_parseval_sqrt_canyon_and_124_theorem_synthesis}.
\begin{enumerate}
  \item \textbf{Entropy Sandwich \& Non-Split Lower Bound (SB.121):} Every $(B, w) \in \nolinkurl{parryCylTree}\ W$ satisfies $1 \le B \le 2^W$, yielding $3^W \le 2^W |\mathcal{C}_W| \le 3^{W+1} - 2 \cdot 2^W$ and $3^k \le 2^k(U_k + U_{k+1})$ for $U_k = 2|\mathcal{C}_k| - |\mathcal{C}_{k+1}|$.
  \item \textbf{Half-Period Antiphase \& Indicator Identity (SB.122):} Splitting children satisfy $((2m+1) b_{\text{right}}) \bmod 2^N = ((2m+1) b_{\text{left}} + 2^{N-1}) \bmod 2^N$, and any $\{0,1\}$-function $f$ on $[0, 2H)$ obeys $\sum_{y < H} (f(y) - f(y+H))^2 + 2 \sum_{y < H} f(y) f(y+H) = \sum_{y < H} (f(y) + f(y+H))$.
  \item \textbf{Conditional $17:12$ Crossover (SB.123):} Since $3^{17} < 2^{27}$, if $2^{5k} S_{\text{card}} \le 2 |\mathcal{C}_{17k}| + 2^{5k} E$ and $2^{17k} E^2 \le C_{\text{err}}^2 3^{17k}$, then $S_{\text{card}} \le (6 + C_{\text{err}}) 2^{5k}$ ($O(X^{5/12})$).
  \item \textbf{Parseval Floor Inequalities (SB.124):} For every $k \in \mathbb{N}$, $3^k \le 2^{k+1} U_k \vee 3^k \le 2^{k+1} U_{k+1}$, and for $W \ge 2$, $3^{3W} < 2^{5W}$ while $2^{4W-3} < 3^{3W-2}$.
\end{enumerate}
\end{theorem}

\begin{theorem}[Coset Split, 3-Point Scaling \& Gap Law (SB.125--SB.132)]
\emph{Lean identifiers:} \nolinkurl{lte2_order_of_three_and_pm_cyclic_partition_of_odd_residues}, \nolinkurl{two_stage_parry_concatenation_cocycle_and_frequency_rotation}, \nolinkurl{conjugate_symmetric_coset_split_and_exact_orbit_mean_identity}, \nolinkurl{three_point_window_scaling_decay_vs_two_point_divergence}, \nolinkurl{zero_three_point_window_triples_implies_at_most_one_positive_survivor}, \nolinkurl{skew_product_strict_growth_and_poincare_anti_recurrence_law}, \nolinkurl{quadratic_abc_two_superblock_mutual_run_cutoff}, \nolinkurl{grand_paths_a_and_b_132_theorem_semigroup_breakout_synthesis}.
\begin{enumerate}
  \item \textbf{$\text{LTE}_2$ Order $2^{K-2}$ \& Coset Half-Sum (SB.125--SB.127):} $3^{2^{m+1}} \bmod 2^{m+4} = 2^{m+3} + 1$ and $3^{2^{m+2}} \bmod 2^{m+4} = 1$, $+\langle 3 \rangle \cap -\langle 3 \rangle = \emptyset$, $(3^{N_1} b_{\text{total}}) \bmod 2^{N_1+N_2} = (3^{N_1} b_{\text{pref}} + 2^{N_1} b_{\text{cont}}) \bmod 2^{N_1+N_2}$, and $(C_{\text{pos}} \cup C_{\text{neg}}).\text{sum } M_{\text{sq}} = 2 \cdot C_{\text{pos}}.\text{sum } M_{\text{sq}}$ under $M_{\text{sq}}(2^K - x) = M_{\text{sq}}(x)$.
  \item \textbf{3-Point Cardinality Scaling \& Zero-Triple Uniqueness (SB.128--SB.129):} $8^N < 9^N \le 4^N |\mathcal{C}_N|^2$, $8^N |\mathcal{C}_N|^3 < 27 \cdot 27^N$, $27 \cdot 4^W \cdot 27^{9W} < 32^{9W}$ ($W \ge 24$), and $(0 \in S \wedge \neg\exists x < y < z \in S) \implies |S \setminus \{0\}| \le 1$.
  \item \textbf{Strict Growth \& Shared-Word Gap (SB.130):} For $b_0 \ge 1$ and $n_1 < n_2$, sharing $L$ parity bits forces $2^L \le \operatorname{orbit}(b_0, n_2) - \operatorname{orbit}(b_0, n_1) < \operatorname{orbit}(b_0, n_2)$.
  \item \textbf{Conditional Quadratic $abc$ Elimination (SB.131--SB.132):} Assuming $3^{2m_1+3} V_1 \le K_2(6 V_1 V_2)^2$ on $2^{2m_2+3} V_2 + 1 = 3^{2m_1+3} V_1$ yields $2^{4m_2+6} < 36 K_2 V_1^3 3^{2m_1+3}$.
\end{enumerate}
\end{theorem}

\begin{theorem}[Method 1--3 Obstructions \& Mod-320 Even-Gauge Pinch (SB.133--SB.140)]
\emph{Lean identifiers:} \nolinkurl{method_one_pigeonhole_cylinder_excess_obstruction}, \nolinkurl{method_two_abc_generic_radical_excess_identity_and_bound}, \nolinkurl{method_three_markov_shift_average_vs_bad_shift_occupancy}, \nolinkurl{interior_odd_core_coprime_to_six_and_prime_divisor_ge_five}, \nolinkurl{dual_multiplicative_euler_product_to_additive_kempner_sandwich}, \nolinkurl{real_half_interval_short_recovery_ban_two_two_through_two_five}, \nolinkurl{real_coupled_gauge_mod64_and_mod320_unique_even_survivor_class}, \nolinkurl{grand_140_theorem_method2_fusion_and_real_gauge_pinch_synthesis}.
\begin{enumerate}
  \item \textbf{Standalone Obstructions (SB.133--SB.135):} (i) $12 k + 1 < 17^k \le |\mathcal{C}_{7k}|$ ($k \ge 1$); (ii) along $s_0 \ge 2, \forall k < M,\, s_{1+k} \ge 2$, $3^{d_0} \le R_0$, $3^{T_K+d_0} u_0 < 2^{S_0}(u_K R_0)$, and $2^{S_0} < 3^{d_0} \implies 2^{S_0} u_K < u_K R_0$ (a product excess, coinciding with $\operatorname{rad}(u_K R_0)$ only when $\gcd(u_K, R_0) = 1$ and both are squarefree); and (iii) $\sum_{a < M} \sum_{j < L} f((a+j) \bmod M) = L \sum_{x < M} f(x)$.
  \item \textbf{Prime Divisors $\ge 5$ \& Excess Sandwich (SB.136--SB.137):} For $1 \le s_r$, $\gcd(u_{r+1}, 6) = 1$ and $(\operatorname{Prime}(p) \wedge p \mid u_{r+1}) \implies p \ge 5$; for $\forall k < K,\, s_k \ge 2$, $3^{d_0} \prod_{k=1}^{K-1} u_k \le \prod_{k=0}^{K-1}(3^{s_k} u_k + 1) - 3^{T_K} \prod_{k=0}^{K-1} u_k < \frac{1}{5} \cdot 3^{T_K} \prod_{k=1}^{K-1} u_k$.
  \item \textbf{Short-Recovery Ban ($1 \le c \le 5$) \& Mod-320 Pinch (SB.138--SB.140):} For $1 \le D, 2 R_0 < D$, and $1 \le c \le 5$, $\Psi_{2+c} = 13 \cdot 3^{c-1}$ implies $\nolinkurl{parryIterBit}\ R_0\ D\ (2+c) \ne 1$ ($\frac{47}{81}, \frac{133}{243}, \frac{383}{729}, \frac{1117}{2187}, \frac{3287}{6561} > \frac{1}{2}$). Along \texttt{ValidBlockChain b\_0 s u (r + 3)} with $2 \le s_r$, $s_{r+1} = 2$, and $6 \le s_{r+2}$, $W_{r+1} \bmod 64 = 4$, $W_{r+1} \bmod 320 = 4$, and $V \bmod 80 = 1$.
\end{enumerate}
\end{theorem}

\begin{theorem}[Kneading Renewal \& Laplace Criteria (SB.141--SB.148)]
\emph{Lean identifiers:} \nolinkurl{parry_kneading_spine_renewal_decomposition_of_cylinder_counts}, \nolinkurl{kneading_renewal_coalesced_one_zero_zero_and_three_scale_phase_lock}, \nolinkurl{multi_scale_phase_non_concentration_and_three_adic_debt_law}, \nolinkurl{pointwise_spectral_gap_to_unconditional_sub_flp_survivor_exponent}, \nolinkurl{pure_nat_euler_sandwich_for_origin_laplace_kernel}, \nolinkurl{origin_laplace_partition_two_way_mahler_threshold_equivalence}, \nolinkurl{discrete_poisson_kernel_strict_real_part_positivity}, \nolinkurl{grand_148_theorem_moves_1_and_2_kneading_and_laplace_synthesis}.
\begin{enumerate}
  \item \textbf{Kneading Renewal, Phase Syzygy \& Conditional $8:7$ Reduction (SB.141--SB.144):} $|\mathcal{C}_N| = 1 + \sum_{j < N,\, 2^{j+1} < 3 B_j} |\mathcal{C}_{N-1-j}|$; $(9 b_{\text{pref}}(9)) \bmod 2^N = (13 b_{\text{pref}}(3)) \bmod 2^N$ ($N \ge 3$); for $1 \le s < K$, $1 \le y_K < 2^K$, $1 \le y_{\text{sub}} < 2^{K-s}$ with $(3^s y_{\text{sub}}) \bmod 2^{K-s} = y_K \bmod 2^{K-s}$ and $d_K + 3^s d_{\text{sub}} < 2^{K-s}$, $d_K = 3^s d_{\text{sub}}$; and if $2^k S_{\text{card}} \le |\mathcal{C}_{8k}| + 2^k E$ and $5^{8k} E \le C_0 7^{8k}$, then $S_{\text{card}} \le 3 \cdot 13^k + C_0 \cdot 15^k$.
  \item \textbf{Euler Sandwich, One-Sided $4/3$ Criteria \& Poisson Core (SB.145--SB.148):} For $m \ge 2$, $2(m-1)^m < m^m$ and $m^{m-1} < 3(m-1)^{m-1}$. For $L \ge 2$ and $0 \in S$: $(\exists b \in S,\, 1 \le b < L \wedge b \le M_{\max}) \implies 4 L^{M_{\max}} < 3 \cdot \nolinkurl{laplacePartitionNum}\ L\ M_{\max}\ S$, while $(1 \le k \wedge (\forall b \in S \setminus \{0\},\, k L \le b \le M_{\max}) \wedge 3|S \setminus \{0\}| \le 2^k) \implies 3 \cdot \nolinkurl{laplacePartitionNum}\ L\ M_{\max}\ S < 4 L^{M_{\max}}$. For $0 \le r < D$ and $|C| \le D$ in $\mathbb{Z}$, $0 < D^2 - 2 r C + r^2 < 2(D^2 - r C)$.
\end{enumerate}
\end{theorem}

\begin{theorem}[$(N+1)$ Widths, Barrier \#7 \& $\mathbb{R}$ Bridges (SB.149--SB.154)]
\emph{Lean identifiers:} \nolinkurl{parry_cylinder_tree_exact_n_plus_one_width_class_compression}, \nolinkurl{upward_hensel_phase_step_dichotomy_and_barrier7_rogue_reset_law}, \nolinkurl{dyadic_frequency_factorization_and_box_high_power_vanishing}, \nolinkurl{grand_152_theorem_mahler_semigroup_breakout_capstone_synthesis}, \nolinkurl{real_z_number_uniqueness_parry_cylinder_bridge_and_flatto_bound}, \nolinkurl{unconditional_real_companion_mod320_even_gauge_pinch}.
\begin{enumerate}
  \item \textbf{Exact $(N+1)$ Scaled-Width Bijection (SB.149):} Every $B_r = \nolinkurl{kneadingSpineB}\ r$ is odd, and $\{B : (B, w) \in \nolinkurl{parryCylTree}\ N\} = \{2^{N-r} B_r : 0 \le r \le N\}$ has exact cardinality $N + 1$ for every $N \in \mathbb{N}$.
  \item \textbf{Upward Hensel Dichotomy \& Rogue-Reset Uniqueness (SB.150):} For $k \ge 1$, $1 \le y_k < 2^k$, $1 \le y_{\text{next}} < 2^{k+1}$ with $y_{\text{next}} \bmod 2^k = (3 y_k) \bmod 2^k$, either $8 d_k \ge 2^k \vee 8 d_{\text{next}} \ge 2^{k+1}$ or $d_{\text{next}} = 3 d_k$; in $\mathbb{Z}$, $2 R \le 3^A \wedge |3^A d_i - q 2^K| < R$ ($i = 1, 2$) forces $d_1 = d_2$.
  \item \textbf{Dyadic Frequency Factorization (SB.151--SB.152):} For $W \le N$ and $1 \le h < 2^N$, $2^{N-W} \mid h \implies (h \cdot 2^W) \bmod 2^N = 0 \wedge h \bmod 2^N \ne 0$, whereas $\neg(2^{N-W} \mid h) \implies h = 2^v h_{\text{odd}}$ with $v < N - W$ and odd $h_{\text{odd}} < 2^{N-v}$.
  \item \textbf{Full Real-Number $\mathbb{R}$ Bridges (SB.153--SB.154):} Over $\mathbb{R}$: (i) \texttt{ValidCompanionForever b\_0} $\alpha \wedge$ \texttt{ValidCompanionForever b\_0} $\beta \implies \alpha = \beta$~\cite{Mahler1968}; (ii) for any $W \in \mathbb{N}$, $X \le 2^W$, and finite set $Z_{\text{set}} \subset [0, X)$ of real Mahler $Z$-numbers, $|Z_{\text{set}}| \le |\mathcal{C}_W|$ and $2^W |Z_{\text{set}}| + 2 \cdot 2^W \le 3^{W+1}$~\cite{Flatto1992,FLP1995}; and (iii) along \texttt{ValidBlockChain b\_0 s u (r + 3)} admitting \texttt{ValidCompanionUpTo b\_0} $\alpha\ (T_{r+3} + 1)$ with $2 \le s_r$ and $s_{r+1} = 2$, survival over $\mathbb{R}$ automatically forces $6 \le s_{r+2}$, $W_{r+1} \bmod 320 = 4$, and $V \bmod 80 = 1$.
\end{enumerate}
\end{theorem}

\section{Conclusion: Seven Structural Barriers Guarding $Z = \emptyset$}

Mahler's conjecture $Z = \emptyset$ remains open. Across \nolinkurl{SemigroupBreakoutTheorems.lean} and the root library, we isolate seven structural barriers that halt standalone proof methods: (1) adelic $\widehat{\mathbb{Z}} \times \mathbb{R}$ decoupling (Meta-Theorem~\ref{thm:meta-decoupling}, \texttt{SB.104}, \texttt{SB.109}); (2) linear Weil height of moving odd cores in Baker--Matveev~\cite{Matveev2000} (\texttt{SB.67}, \texttt{SB.71}); (3) Ridout's $Q_m^{-2}$ critical exponent~\cite{Ridout1957} (\texttt{SB.75}); (4) Schlickewei's 2D subspace syzygy on $\ge 3$ shifted origins~\cite{Schlickewei1990}; (5) the P\'olya--Carlson capacity threshold $R_\infty R_2 = 1$~\cite{Polya1921} (\texttt{SB.3}, \texttt{SB.117}); (6) the Parseval odd-mode floor and $b_0 = 0$ dyadic floor (\texttt{SB.124}, \texttt{SB.128}--\texttt{SB.129}); and (7) Pigeonhole cylinder excess $17^k > 12k+1$, generic coprime product excess $u_K R_0 > 2^{S_0} u_K$, and variable-$d$ rogue resets (\texttt{SB.133}, \texttt{SB.134}, \texttt{SB.150}).

\section*{Acknowledgments}
Developed by Joshua Coon, with computational exploration, Lean~4 formalization assistance, and adversarial manuscript review from the Astra, Fable, and Argon AI models in GitHub Copilot.

\begin{thebibliography}{99}
\bibitem{Mahler1968} K.~Mahler, ``An unsolved problem on the powers of $3/2$,'' \emph{J.\ Austral.\ Math.\ Soc.}\ \textbf{8} (1968), 313--321.
\bibitem{Ridout1957} D.~Ridout, ``Rational approximations to algebraic numbers,'' \emph{Mathematika}\ \textbf{4} (1957), 125--131.
\bibitem{Polya1921} G.~P\'olya, ``\"Uber gewisse notwendige Determinantenkriterien f\"ur die Fortsetzbarkeit einer Potenzreihe,'' \emph{Math.\ Ann.}\ \textbf{99} (1921), 687--706.
\bibitem{Parry1960} W.~Parry, ``On the $\beta$-expansions of real numbers,'' \emph{Acta Math.\ Acad.\ Sci.\ Hungar.}\ \textbf{11} (1960), 401--416.
\bibitem{SengeStraus1973} H.~G.~Senge and E.~G.~Straus, ``PV-numbers and sets of multiplicity,'' \emph{Period.\ Math.\ Hungar.}\ \textbf{3} (1973), 93--100.
\bibitem{Stewart1980} C.~L.~Stewart, ``On the representation of an integer in two different bases,'' \emph{J.\ Reine Angew.\ Math.}\ \textbf{319} (1980), 63--72.
\bibitem{Schlickewei1990} H.~P.~Schlickewei, ``$S$-unit equations over number fields,'' \emph{Invent.\ Math.}\ \textbf{102} (1990), 95--107.
\bibitem{Flatto1992} L.~Flatto, ``$Z$-numbers and $\beta$-transformations,'' in \emph{Symbolic Dynamics and Its Applications}, Contemp.\ Math.\ \textbf{135}, AMS, 1992, pp.~181--201.
\bibitem{FLP1995} L.~Flatto, J.~C.~Lagarias, and A.~D.~Pollington, ``On the range of fractional parts $\{\xi(p/q)^n\}$,'' \emph{Acta Arith.}\ \textbf{70} (1995), 125--147.
\bibitem{Matveev2000} E.~M.~Matveev, ``An explicit lower bound for a homogeneous rational linear form in the logarithms of algebraic numbers.\ II,'' \emph{Izv.\ Math.}\ \textbf{64} (2000), 1217--1269.
\bibitem{Dubickas2008} A.~Dubickas, ``On the powers of $3/2$ and other rational numbers,'' \emph{Math.\ Nachr.}\ \textbf{281} (2008), 951--958.
\bibitem{Dubickas2009} A.~Dubickas, ``Powers of a rational number modulo $1$ cannot lie in a small interval,'' \emph{Acta Arith.}\ \textbf{137} (2009), 233--239.
\bibitem{BourgainGamburd2008} J.~Bourgain and A.~Gamburd, ``Uniform expansion bounds for Cayley graphs of $\mathrm{SL}_2(\mathbb{F}_p)$,'' \emph{Ann.\ of Math.}\ (2) \textbf{167} (2008), 625--642.
\end{thebibliography}

\end{document}
